\chapter{Conclusões e Recomendações}
\label{cap:05}

Desenvolver um jogo foi uma tarefa complexa que demanda esforços não apenas lógicos mas também criativos. Ao longo do projeto, foi possível aplicar na prática diversos conhecimentos adquiridos durante o curso, como lógica de programação, estruturas de dados e organização de código. Além disso, foi necessário aprender e se adaptar a novas ferramentas, como a \textit{engine Godot}, se mostrou bastante completa e funcional para o desenvolvimento 2D.

Durante o desenvolvimento, diversos desafios foram enfrentados, especialmente em relação com aprendizado da \textit{engine} e toda a estrutura que ela traz. Entender como organizar os elementos do jogo, configurar o comportamento do personagem e gerenciar as animações exigiu bastante estudo e prática. Esses obstáculos, no entanto, contribuíram para uma melhor compreensão de conceitos importantes e reforçaram a importância de uma base sólida em lógica e estruturação de código.

A experiência de criar um jogo também trouxe uma percepção mais clara sobre a importância do planejamento e da organização ao longo do processo. Foi necessário dividir as tarefas, priorizar funcionalidades e realizar testes constantes para garantir o funcionamento correto das mecânicas. Com isso, ficou evidente como o desenvolvimento de jogos é um processo iterativo que exige paciência, atenção aos detalhes e capacidade de adaptação.

Por fim, o projeto representou uma oportunidade significativa para aplicar o conhecimento teórico em um contexto prático e criativo. A finalização de um jogo funcional e jogável, após diversas etapas de desenvolvimento, evidenciou a eficácia do processo e das ferramentas utilizadas. Além disso, o trabalho reforçou o potencial da área de desenvolvimento de jogos, que integra programação, design e criatividade em um objetivo comum.
